\TieudeT{PHONG TỎA VẬT LÍ}
\subsubsection*{PHẦN I. Thí sinh trả lời câu hỏi từ câu 1 đến câu 18. Mỗi câu hỏi thí sinh chỉ chọn một phương án}
\setcounter{ex}{0}
%mac
%\loai{Biểu thức}
\begin{ex}%book%[1P3Y8-1][Loai1]
Một ngẫu lực gồm hai lực $\vec F_1 $ và $\vec F_2 $ có $F_1 = F_2 = F$ và có cánh tay đòn $d$. Moment của ngẫu lực này là
\choice
{\True $Fd$}
{$\left(F_1 - F_2\right).d$}
{$\left(F_1 + F_2\right).d$}
{Chưa đủ dữ kiện để tính toán}
\loigiai{  
Mômen ngẫu lực: $M = F.d$.
}
\end{ex}
\begin{ex}%book%[1P3Y8-1][Loai1]
Moment $M$ của ngẫu lực được tính bằng công thức
\choice
{$M = Fd$ với $F$ là tổng độ lớn hai lực và $d$ là khoảng cách giữa hai giá của hai lực}
{$M=\dfrac{F}{d}$ với $F$ là độ lớn của một lực và $d$ là khoảng cách giữa hai giá của hai lực}
{\True $M = F.d$ với $F$ là độ lớn của một lực và $d$ là khoảng cách giữa hai giá của hai lực}
{$M = F.d$ với $F$ là độ lớn của một lực và $d$ là đoạn nối hai điểm đặt của hai lực}
\loigiai{
Ta có: $M = Fd$. Với F là độ lớn của một lực và d là khoảng cách giữa hai giá của hai lực.
}
\end{ex}
%\loai{Đơn vị}
\begin{ex}%book%[1P3Y8-1][Loai1]
Đơn vị của moment ngẫu lực là
\choice
{$N$}
{$W$}
{\True $N.m$}
{$J$}
\loigiai{	      
Do $M = F.d$ với $[F] = N, [d] = m$ $\Rightarrow$ $[M] = N.m$.
}
\end{ex}

%\loai{Khái niệm}
\begin{ex}%book%[1P3Y8-1][Loai1]
Ngẫu lực là
\choice
{hợp lực của hai lực song song cùng chiều}
{hợp của hai lực song song ngược chiều}
{\True hệ hai lực song song, ngược chiều, có độ lớn bằng nhau cùng tác dụng lên một vật}
{hệ hai lực song song, cùng chiều, có độ lớn bằng nhau cùng tác dụng lên một vật}
\loigiai{      
Hệ hai lực song song, ngược chiều, có độ lớn bằng nhau và cùng tác dụng vào một vật.
}
\end{ex}
%\loai{Đặc điểm}
\begin{ex}%book%[1P3Y8-1][Loai1]
Phát biểu nào sau đây đúng khi nói về ngẫu lực tác dụng lên một vật?
\choice
{Ngẫu lực là hệ hai lực cùng giá, ngược chiều, cùng độ lớn}
{\True Ngẫu lực là hệ hai lực song song, ngược chiều và có cùng độ lớn}
{Ngẫu lực là cặp lực có thể tổng hợp thành một lực duy nhất}
{Ngẫu lực làm cho vật chuyển động tròn đều}
\loigiai{	      
Ngẫu lực là hệ hai lực song song, ngược chiều, có cùng độ lớn. Ngẫu lực có tác dụng làm vật quay.
}
\end{ex}
\begin{ex}%book%[1P3Y8-1][Loai1]
Ngẫu lực có tác dụng
\choice
{làm vật rắn chuyển động tịnh tiến}
{\True làm vật rắn chuyển động quay}
{làm vật rắn vừa chuyển động tịnh tiến vừa chuyển động quay}
{làm vật rắn quay quanh trục bất kì đi qua trọng tâm}
\loigiai{
Ngẫu lực có tác dụng làm vật rắn quay.
}
\end{ex}
\begin{ex}%book%[1P3Y8-1][Loai1]
Hai lực tác dụng vào một vật tạo thành ngẫu lực khi
\choice
{cùng trên một giá, ngược chiều và có độ lớn bằng nhau}
{song song, cùng chiều và có độ lớn không bằng nhau}
{song song, ngược chiều và có độ lớn không bằng nhau}
{\True song song, ngược chiều và có độ lớn bằng nhau}
\loigiai{
Song song, ngược chiều và có độ lớn bằng nhau là đáp án đúng.
}
\end{ex}
\begin{ex}%Tailieuchuan
\immini[thm]{Khi đạp xe, moment ngẫu lực nào dưới đây sẽ giúp xe đạp tăng tốc nhanh nhất?
\choice
{Moment lớn, tác dụng ngắt quãng lên bàn đạp}
{Moment nhỏ, tác dụng đều và liên tục lên bàn đạp}
{\True Moment lớn, tác dụng đều và liên tục lên bàn đạp}
{Moment nhỏ, tác dụng ngắt quãng lên bàn đạp}}{
\includegraphics[width=4cm]{images/2025_04_24_a80008410be527bbd242g-13(1)}}
\loigiai{}
\end{ex}
\begin{ex}%Tailieuchuan
Cặp lực nào trong hình bên dưới là ngẫu lực?
\choice
{\begin{tikzpicture}[scale=1,thick,>=stealth',draw=Mapcolor]
\path
(0,0)coordinate(O)
(0.5,0)coordinate(A)
(-0.5,0)coordinate(B)
;
\draw  plot[smooth, tension=.7] coordinates {(-0.5,-1) (-1.5,-1) (-1.5,-0.5) (-1.5,0.5) (-0.5,1) (0.5,1) (1.5,0.5) (1.5,-0.5) (1,-1) (-0.5,-1)};
\draw[dashed] (A)--(B);
\draw[->,Mapcolor] (-0.5,0) -- (-2,0);
\draw[->,Mapcolor] (0.5,0) -- (2,0);
\node at (-2,-0.5) {$\vec F_1$};
\node at (2,-0.5) {$\vec F_2$};
\foreach \x/\g in {A/90,B/180}
\draw[fill=white](\x) circle(0.085)+(\g:0.3)node{};
\end{tikzpicture}}
{\True \begin{tikzpicture}[scale=1,thick,>=stealth',draw=Mapcolor]
\path
(0,0)coordinate(O)
(1.5,-0.5)coordinate(A)
(-0.5,0.5)coordinate(B)
;
\draw  plot[smooth, tension=.7] coordinates {(-0.5,-1) (-1.5,-1) (-1.5,-0.5) (-1.5,0.5) (-0.5,1) (0.5,1) (2.5,1) (2.5,-1) (1,-1) (-0.5,-1)};
\draw[->,Mapcolor] (-0.5,0.5) -- (1,0.5);
\draw[->,Mapcolor] (1.5,-0.5) -- (0,-0.5);
\node at (1.5,0.5) {$\vec F_1$};
\node at (-0.5,-0.5) {$\vec F_2$};
\foreach \x/\g in {A/90,B/180}
\draw[fill=white](\x) circle(0.085)+(\g:0.3)node{};
\end{tikzpicture}}
{\begin{tikzpicture}[scale=1,thick,>=stealth',draw=Mapcolor]
\path
(0,0)coordinate(O)
(0.5,-0.5)coordinate(A)
(0.5,0.5)coordinate(B)
;
\draw  plot[smooth, tension=.7] coordinates {(-0.5,-1) (-1.5,-1) (-1.5,-0.5) (-1.5,0.5) (-0.5,1) (0.5,1) (2.5,0.5) (2.5,-1) (1,-1) (-0.5,-1)};
\draw[->,Mapcolor] (0.5,0.5) -- (3.5,0.5);
\draw[->,Mapcolor] (0.5,-0.5) -- (1.5,-0.5);
\node at (4,0.5) {$\vec F_1$};
\node at (2,-0.5) {$\vec F_2$};
\foreach \x/\g in {A/90,B/180}
\draw[fill=white](\x) circle(0.085)+(\g:0.3)node{};
\end{tikzpicture}
}
{\begin{tikzpicture}[scale=1,thick,>=stealth',draw=Mapcolor]
\path
(0,0)coordinate(O)
(0.5,-0.5)coordinate(A)
(0.5,0.5)coordinate(B)
;
\draw  plot[smooth, tension=.7] coordinates {(-0.5,-1) (-1.5,-1) (-1.5,-0.5) (-1.5,0.5) (-0.5,1) (0.5,1) (2.5,0.5) (2.5,-1) (1,-1) (-0.5,-1)};
\draw[->,Mapcolor] (0.5,0.5) -- (3.5,0.5);
\draw[->,Mapcolor] (0.5,-0.5) -- (1.5,-0.5);
\node at (4,0.5) {$\vec F_2$};
\node at (2,-0.5) {$\vec F_1$};
\foreach \x/\g in {A/90,B/180}
\draw[fill=white](\x) circle(0.085)+(\g:0.3)node{};
\end{tikzpicture}}

\loigiai{

}
\end{ex}
\begin{ex}%Tailieuchuan
Nếu moment ngẫu lực tác dụng lên một vật đang quay quanh trục bằng $0$ thì vật sẽ
\choice
{\True quay đều}
{quay nhanh dần}
{không quay}
{chuyển động thẳng đều}
\loigiai{}
\end{ex}
\begin{ex}%Tailieuchuan
Một vật không có trục quay cố định, khi chịu tác dụng của một ngẫu lực thì vật sẽ
\choice
{chuyển động tịnh tiến}
{\True chuyển động quay}
{vừa quay, vừa tịnh tiến}
{nằm cân bằng}
\loigiai{

}
\end{ex}
\begin{ex}%Tailieuchuan
Hai lực của một ngẫu lực có độ lớn $F$. Cánh tay đòn của ngẫu lực $d$. Muốn moment ngẫu lực tăng gấp đôi và giảm $F$ đi một nửa thì $d$ phải tăng lên
\choice
{$1$ lần}
{\True $4$ lần}
{$8$ lần}
{$2$ lần}
\loigiai{}
\end{ex}

\begin{ex}%Tailieuchuan
Khi dùng tua vít để vặn đinh vít, người ta đã tác dụng vào các đinh vít
\choice
{\True một ngẫu lực}
{hai ngẫu lực}
{cặp lực cân bằng}
{cặp lực trực đối}
\loigiai{}
\end{ex}

\begin{ex}%Tailieuchuan
\immini[thm]{Moment ngẫu lực đối với trục quay $O$ vuông góc với mặt phẳng của ngẫu lực như hình vẽ. Công thức nào sau đây đúng?
\choice
{\True $M=F_1 \ d_1+F_2 \ d_2$}
{$M=\left|F_1 \ d_1-F_2 \ d_2\right|$}
{$M=F_1 \ d_2+F_2 \ d_1$}
{$M=\left|F_1 \ d_2-F_2 \ d_1\right|$}}{
\begin{tikzpicture}[scale=1.25,thick,>=stealth',draw=Mapcolor]
\path
(0,0)coordinate(O)
(0,-0.5)coordinate(A)
(0,0.5)coordinate(B)
;
\draw  plot[smooth, tension=.7] coordinates {(-0.5,-1) (-1.5,-1) (-1.5,-0.5) (-1.5,0.5) (-0.5,1) (0.5,1) (1.5,0.5) (1.5,-0.5) (1,-1) (-0.5,-1)};
\draw[->,Mapcolor] (A)--(B);
\draw[->,Mapcolor] (0,0.5) -- (2.5,0.5);
\draw[->,Mapcolor] (0,-0.5) -- (-2.5,-0.5);
\node at (3,0.5) {$\vec F_1$};
\node at (-2.5,-1) {$\vec F_2$};
\foreach \x/\g in {A/90,B/180,O/90}
\draw[fill=white](\x) circle(0.05)+(\g:0.3)node{};
\node at (0.2,0.3) {$d_1$};
\node at (0.2,-0.2) {$d_2$};
\end{tikzpicture}}
\loigiai{}
\end{ex}
\begin{ex}%Tailieuchuan
Hai lực của một ngẫu lực có độ lớn $F=20 \ N$. Cánh tay đòn của ngẫu lực $d=30 \ cm$. Moment của ngẫu lực là
\choice
{$600\ N.m$}
{$60\ N.m$}
{\True $6\ N.m$}
{$0,6\ N.m$}
\loigiai{}
\end{ex}

\begin{ex}%Tailieuchuan
\immini[thm]{Một người tài xế dùng hai tay để quay tay lái ô tô. Hai lực có độ lớn bằng nhau là $F=40 \ N$, tác dụng ngược chiều và song song vào hai điểm trên vành tay lái, cách nhau một khoảng $d=0,3 \ m$. Moment ngẫu lực tác dụng lên tay lái có độ lớn bằng
\choice
{\True $12\ N.m$}
{$10\ N.m$}
{$15\ N.m$}
{$24\ N.m$}}{
\includegraphics[width=4cm]{images/2025_04_24_a80008410be527bbd242g-09(1)}}
\loigiai{}
\end{ex}










%mac
%\loai{Tính moment ngẫu lực}
\begin{ex}%book%[1P3Y8-2][Loai1]
Hai lực của một ngẫu lực có độ lớn $F = 5,0\ N$. Cánh tay đòn của ngẫu lực $d = 20\ cm$. Moment của ngẫu lực là
\choice
{$100\ N.m$}
{$2\ N.m$}
{$0,5\ N.m$}
{\True $1,0\ N.m$}
\haidong{3}		\loigiai{Moment của ngẫu lực là: $M = Fd = 5.0,2 = 1$ N.m}
\end{ex}
\begin{ex}%book%[1P3Y8-2][Loai1]
Hai lực của ngẫu lực có độ lớn $F = 20\ N$. Cánh tay đòn của ngẫu lực $d = 30\ cm$. Moment của ngẫu lực là
\choice
{$600\ N.m$}
{$60\ N.m$}
{\True $6\ N.m$}
{$0,6\ N.m$}
\haidong{3}		\loigiai{
Ta có: $M = F.d = 20.0,3 = 6\ N.m$.
}
\end{ex}
%\loai{Tính moment - hình tam giác}

\subsubsection*{PHẦN 2. Thí sinh trả lời từ câu 1 đến câu 4. Trong mỗi ý a), b), c), d) ở mỗi câu, thí sinh chọn đúng hoặc sai.}
\setcounter{ex}{0}
\begin{ex}
\immini[thm]{
Cho hai lực $\vec{F}_1$ và $\vec{F}_2$ tác dụng vào vật như hình vẽ
\choiceTF[t]
{Hai lực $\vec{F}_1$ và $\vec{F}_2$ làm cho vật quay theo hai chiều ngược nhau}
{\True Hai lực $\vec{F}_1$ và $\vec{F}_2$ được gọi là ngẫu lực}
{\True $d=d_1+d_2$ được gọi là cánh tay đòn của ngẫu lực}
{Moment của ngẫu lực phụ thuộc vào vị trí trục quay vuông góc với mặt phẳng chứa ngẫu lực}}{\begin{tikzpicture}[scale=1,thick,>=stealth',draw=Mapcolor]
\path
(0,0)coordinate(O)
;

\draw[fill=Mapcolor!30]  (0,0) rectangle (0.5,-5.5);
\draw[dashed] (-1.5,0) -- (0,0);

\draw[dashed] (-1.5,-4) -- (1.75,-4);

\draw[dashed] (0.5,-5.5) -- (2.25,-5.5);
\draw[<->] (1,0) -- (1,-4);
\draw [<->](1,-4) -- (1,-5.5);
\draw[<->] (2,0) -- (2,-5.5);
\draw[->,very thick,blue] (0.25,0) -- (2.5,0);
\draw[->,very thick,blue] (0.25,-5.5) -- (-2,-5.5);
\node at (1.25,-2) {$d_1$};
\node at (1.25,-4.75) {$d_2$};
\node at (3.25,-2.5) {$d=d_1+d_2$};
\end{tikzpicture}}
\loigiai{
\begin{itemchoice}
\itemch
\itemch
\itemch
\itemch
\end{itemchoice}}
\end{ex}
\begin{ex}
\immini[thm]{
Tại hai điểm $A$ và $B$ của một vật cách nhau $AB=40\ cm$, người ta tác dụng hai lực song song, ngược chiều và có độ lớn $F_1=F_2=0,49\ N$ như hình vẽ.}{\begin{tikzpicture}[scale=1,thick,>=stealth',draw=Mapcolor]
\path
(0,0)coordinate(O)
;
\draw [fill=Mapcolor!30] (0,0.25) rectangle (5.75,0);
\draw[very thick,->,Mapcolor] (0,0) -- (1.25,-1.25);
\draw[very thick,->,Mapcolor] (5.75,0.25) -- (4.5,1.25);
\node at (0,0.5) {A};
\node at (5.75,0.5) {B};
\node at (1.5,-1.5) {$\vec F_1$};
\node at (4.25,1.5) {$\vec F_2$};
\node at (0.75,-0.25) {$30^\circ$ };
\end{tikzpicture}}
\choiceTF[t]
{\True Hai lực $\vec{F}_1$ và $\vec{F}_2$ được gọi là ngẫu lực}
{Moment của ngẫu lực phụ thuộc vào vị trí trục quay vuông góc với mặt phẳng chứa ngẫu lực}
{Cánh tay đòn của ngẫu lực là $AB=40\ cm$}
{\True Moment ngẫu lực là $0,098\ N.m$}
\loigiai{
\begin{itemchoice}
\itemch
\itemch
\itemch
\itemch
\end{itemchoice}}
\end{ex}
\begin{ex}
\immini[thm]{Trên một ổ khóa cửa một cánh cửa có hình quả đấm, người ta tác dụng một ngẫu lực $\vec{F}$ được mô tả như hình vẽ. Cho biết $OA=5\ cm$; $F=10\ N$; $\alpha=30^\circ$.}{
\begin{tikzpicture}[scale=.75,thick,>=stealth',draw=Mapcolor]
\path
(0,0)coordinate(O)
;
\draw (O) circle(2cm);
\draw[dashed] (-1.5,1.25) -- (1.5,-1.25);
\draw[very thick,->,Mapcolor] (-1.5,1.25) -- (-3.25,1.25);
\draw [very thick,->,Mapcolor] (1.25,-1.25) -- (3.25,-1.25);
\draw[dashed] (-1.5,1.25) -- (2.5,1.25);
\node at (-1.5,1.75) {A};
\node at (1.5,-1.75) {B};
\node at (-3.25,0.75) {$\vec F_1$};
\node at (3.25,-1.75) {$\vec F_2$};
\node at (-0.75,1) {$\alpha$};
\foreach \x/\g in {O/90}
\draw[fill=white](\x) circle(0.05)+(\g:0.3)node{$\x$};
\end{tikzpicture}}
\choiceTF[t]
{\True Chiều quay của ổ khóa theo chiều quay của kim đồng hồ}
{Cánh tay đòn của mỗi lực $\vec{F}$ bằng $5\ cm$}
{\True Cánh tay đòn của ngẫu lực bằng $5\ cm$}
{\True Momen ngẫu lực bằng $0,5\ N.m$}
\loigiai{
\begin{itemchoice}
\itemch
\itemch
\itemch
\itemch
\end{itemchoice}}
\end{ex}
\begin{ex}%Tailieuchuan
Một vật chịu tác dụng của một ngẫu lực gồm hai lực song song, ngược chiều, có độ lớn bằng nhau.
\choiceTF[t]
{Moment ngẫu lực chỉ phụ thuộc vào khoảng cách từ giá của một trong hai lực đến trục quay} 
{\True Hai lực trong ngẫu lực cùng có tác dụng làm vật quay theo một chiều} 
{\True Đơn vị đo của moment ngẫu lực là $N.m$} 
{Nếu mỗi lực có độ lớn là $5\ N$ và khoảng cách hai giá của hệ ngẫu lực bằng $1\ m$ thì moment ngẫu lực có độ lớn $10\ N.m$} 
\loigiai{ 
\begin{itemchoice}
\itemch
\itemch
\itemch
\itemch
\end{itemchoice}
}
\end{ex}






















\subsubsection*{PHẦN 3. Thí sinh trả lời từ câu 1 đến câu 6.}
\setcounter{ex}{0} 
\begin{ex}%Tailieuchuan
Mỗi lực của một ngẫu lực có độ lớn $F=40 \ N$. Cánh tay đòn của ngẫu lực là $d=30 \ cm$. Moment của ngẫu lực là bao nhiêu $N.m$?
\shortans[oly]{12}
\loigiai{}
\end{ex}

\begin{ex}%Tailieuchuan
Mỗi lực của một ngẫu lực có độ lớn $F=5,0\ N$. Cánh tay đòn của ngẫu lực $d=20\ cm$. Moment của ngẫu lực bằng bao nhiêu N.m?
\shortans[oly]{1}
\loigiai{
}
\end{ex}
\begin{ex}%Tailieuchuan
Mỗi lực của một ngẫu lực có độ lớn $10\ N$, cánh tay đòn của ngẫu lực bằng bao nhiêu $cm$ nếu moment ngẫu lực có độ lớn $2\ N.m$?
\shortans[oly]{20}
\loigiai{}
\end{ex}
\begin{ex}%book%[1P3B8-2][Loai1]
Một vật rắn phẳng, mỏng có dạng là một tam giác đều $ABC$, mỗi cạnh là $a = 20\ cm$. Người ta tác dụng vào vật một ngẫu lực nằm trong mặt phẳng của tam giác. Các lực có độ lớn là $8,0\ N$ đặt vào hai đỉnh $A$ và $B$. Moment của ngẫu lực khi các lực vuông góc với cạnh $AB$ có độ bao nhiêu $N.m$?
\shortans[oly]{1,6}
\haidong{3}		\loigiai{
Ta có: $M = F.d = F.a = 1,6\ N.m$.
}
\end{ex}
\begin{ex}%Tailieuchuan
Một vật rắn chịu tác dụng của lực $F$ quay quanh một trục, khoảng cách từ giá của lực đến trục quay là $d$. Khi tăng lực tác dụng lên  $6$ lần và giảm $d$ đi $2$ lần thì moment của lực $F$ tác dụng lên vật tăng lên bao nhiêu lần?
\shortans[oly]{3}
\loigiai{}
\end{ex}
\begin{ex}
Một vật rắn phẳng mỏng có dạng là một hình vuông $ABCD$, mỗi cạnh là $a=10\ cm$. Người ta tác dụng một ngẫu lực nằm trong mặt phẳng của hình vuông. Biết các lực vuông góc với đường chéo $AC$ có độ lớn $10\ N$ và đặt vào hai đỉnh $A$ và $C$. Moment của ngẫu lực là bao nhiêu $N.m$ (làm tròn kết quả đến chữ số hàng phần trăm)?
\shortans[oly]{1,41}
\loigiai{}
\end{ex}




